%-------------------------
% Resume in Latex
% Author : Jake Gutierrez (template), filled in for Jasper Zhang
% Based off of: https://github.com/sb2nov/resume
% License : MIT
%------------------------

\documentclass[letterpaper,11pt]{article}

\usepackage{latexsym}
\usepackage[empty]{fullpage}
\usepackage{titlesec}
\usepackage{marvosym}
\usepackage[usenames,dvipsnames]{color}
\usepackage{verbatim}
\usepackage{enumitem}
\usepackage[hidelinks]{hyperref}
\usepackage{fancyhdr}
\usepackage[english]{babel}
\usepackage{tabularx}
\input{glyphtounicode}


%----------FONT OPTIONS----------
% sans-serif
% \usepackage[sfdefault]{FiraSans}
% \usepackage[sfdefault]{roboto}
% \usepackage[sfdefault]{noto-sans}
% \usepackage[default]{sourcesanspro}

% serif
% \usepackage{CormorantGaramond}
% \usepackage{charter}


\pagestyle{fancy}
\fancyhf{} % clear all header and footer fields
\fancyfoot{}
\renewcommand{\headrulewidth}{0pt}
\renewcommand{\footrulewidth}{0pt}

% Adjust margins
\addtolength{\oddsidemargin}{-0.5in}
\addtolength{\evensidemargin}{-0.5in}
\addtolength{\textwidth}{1in}
\addtolength{\topmargin}{-.5in}
\addtolength{\textheight}{1.0in}

\urlstyle{same}

\raggedbottom
\raggedright
\setlength{\tabcolsep}{0in}

% Sections formatting
\titleformat{\section}{
  \vspace{-4pt}\scshape\raggedright\large
}{}{0em}{}[\color{black}\titlerule \vspace{-5pt}]

% Ensure that generate pdf is machine readable/ATS parsable
\pdfgentounicode=1

%-------------------------
% Custom commands
\newcommand{\resumeItem}[1]{
  \item\small{
    {#1 \vspace{-2pt}}
  }
}

\newcommand{\resumeSubheading}[4]{
  \vspace{-2pt}\item
    \begin{tabular*}{0.97\textwidth}[t]{l@{\extracolsep{\fill}}r}
      \textbf{#1} & #2 \\
      \textit{\small#3} & \textit{\small #4} \\
    \end{tabular*}\vspace{-7pt}
}

\newcommand{\resumeSubSubheading}[2]{
    \item
    \begin{tabular*}{0.97\textwidth}{l@{\extracolsep{\fill}}r}
      \textit{\small#1} & \textit{\small #2} \\
    \end{tabular*}\vspace{-7pt}
}

\newcommand{\resumeProjectHeading}[2]{
    \item
    \begin{tabular*}{0.97\textwidth}{l@{\extracolsep{\fill}}r}
      \small#1 & #2 \\
    \end{tabular*}\vspace{-7pt}
}

\newcommand{\resumeSubItem}[1]{\resumeItem{#1}\vspace{-4pt}}

\renewcommand\labelitemii{$\vcenter{\hbox{\tiny$\bullet$}}$}

\newcommand{\resumeSubHeadingListStart}{\begin{itemize}[leftmargin=0.15in, label={}]}
\newcommand{\resumeSubHeadingListEnd}{\end{itemize}}
\newcommand{\resumeItemListStart}{\begin{itemize}}
\newcommand{\resumeItemListEnd}{\end{itemize}\vspace{-5pt}}

%-------------------------------------------
%%%%%%  RESUME STARTS HERE  %%%%%%%%%%%%%%%%%%%%%%%%%%%%


\begin{document}

%----------HEADING----------
\begin{center}
    \textbf{\Huge \scshape Jasper Zhang} \\ \vspace{1pt}
    \small 516-288-1886 $|$ \href{mailto:jasperzhang1886@gmail.com}{\underline{jasperzhang1886@gmail.com}} $|$
    \href{https://linkedin.com/in/jasperzhang1001}{\underline{linkedin.com/in/jasperzhang1001}} $|$
    \href{https://github.com/JasperZG}{\underline{github.com/JasperZG}}
\end{center}


%-----------EDUCATION-----------
\section{Education}
  \resumeSubHeadingListStart
    \resumeSubheading
      {William A. Shine Great Neck South High School}{Great Neck, NY}
      {Class of 2027 $|$ GPA: 100.5/100 (weighted)}{Sep. 2023 -- Jun. 2027}
      \resumeItemListStart
        \resumeItem{\textbf{Coursework:} AP Calculus BC, AP Chemistry, AP Physics 2, AP US History, AP English Language.}
        \resumeItem{\textbf{Honors:} U.S. Earth Science Olympiad Top 45 National Finalist (2026); U.S. Presidential AI Challenge New York State Champion (2026); ASM International Special Research Award (2025); AIME Qualifier (2025); Long Island Young Mathematics Scholar Award (2024)}
      \resumeItemListEnd
  \resumeSubHeadingListEnd


%-----------EXPERIENCE-----------
\section{Experience}
  \resumeSubHeadingListStart

    \resumeSubheading
      {Machine Learning \& Computational Biology Researcher}{2025 -- Present}
      {Great Neck South Research Facility}{Great Neck, NY}
      \resumeItemListStart
        \resumeItem{Lead first/second author across six research papers in machine learning for biology, working on problems spanning mechanistic interpretability, multi-task learning theory, and molecular and genomic property prediction; accepted to ICLR 2026 workshops with 16 total paper acceptances and 2 oral presentations (ICLR GEM Workshop, LION20).}
      \resumeItemListEnd

    \resumeSubheading
      {Electrochemistry Researcher}{Summer 2025}
      {Tufts University --- Leverick Group}{Medford, MA}
      \resumeItemListStart
        \resumeItem{Led an independent lithium-ion battery project under Prof. Leverick, designing cell chemistries and fabricating full cells from electrode preparation through assembly; applied scanning electron microscopy (SEM) and related characterization techniques to analyze electrode/electrolyte microstructure and post-cycling degradation.}
      \resumeItemListEnd

  \resumeSubheading
      {Biochemistry Research Assistant}{Summer 2024}
      {Stony Brook University --- Tonge Group}{Stony Brook, NY}
      \resumeItemListStart
        \resumeItem{Contributed to ongoing development of small-molecule inhibitors targeting enzyme drug targets implicated in cancer pathways; cultured E. coli for recombinant protein expression and gene extraction, performing spectrophotometry, plasmid isolation, column chromatography, and culture incubation.}
      \resumeItemListEnd

  \resumeSubHeadingListEnd


%-----------PROJECTS-----------
\section{Projects}
\small{\textit{$^{*}$denotes co-first authorship.}}\vspace{4pt}

\begin{itemize}[leftmargin=0.15in, label={}, itemsep=5pt, topsep=2pt]
    \small
    \item \textbf{The Mechanistic Invariance Test: Genomic Language Models Fail To Learn Positional Regulatory Logic.} J. Zhang$^{*}$, B. Cheng$^{*}$. \emph{ICLR 2026 GEM Workshop (Oral)}.
        \begin{itemize}[leftmargin=0.2in, itemsep=0pt, topsep=2pt]
            \item Built a benchmark probing whether genomic language models learn positional regulatory logic; evaluated five major architectures spanning autoregressive, masked, and state-space paradigms, and showed that apparent positional sensitivity is driven by AT-content shortcuts, with a 100-parameter baseline outperforming billion-parameter models.
        \end{itemize}

    \item \textbf{Single-Position Intervention Fails: Distributed Output Templates Drive In-Context Learning.} \\
    J. Zhang$^{*}$, B. Cheng$^{*}$. \emph{LION20 (Oral)}.
        \begin{itemize}[leftmargin=0.2in, itemsep=0pt, topsep=2pt]
            \item Used causal activation interventions to show that in-context task identity is distributed across demonstration tokens rather than localized; single-position edits achieve 0\% transfer while multi-position edits reach up to 96\% across LLaMA, Qwen, and Gemma.
        \end{itemize}

    \item \textbf{Information-Theoretic Requirements for Gradient-Based Task Affinity Estimation in Multi-Task Learning.} \\
    J. Zhang$^{*}$, B. Cheng$^{*}$. \emph{ICLR 2026 GEM Workshop}.
        \begin{itemize}[leftmargin=0.2in, itemsep=0pt, topsep=2pt]
            \item Identified a sharp phase transition in gradient-based task affinity estimation around $\sim$30--40\% sample overlap; showed that benchmarks like MoleculeNet and TDC operate far below this threshold, explaining years of inconsistent multi-task learning results.
        \end{itemize}

\end{itemize}

%-----------TECHNICAL SKILLS-----------
\section{Technical Skills}
 \begin{itemize}[leftmargin=0.15in, label={}]
    \small{\item{
     \textbf{Programming \& ML}{: Python, Java; PyTorch, NumPy, pandas, scikit-learn, Matplotlib, LaTex; deep learning, GANs, CNNs, language models, multi-task learning, mechanistic interpretability} \\
     \textbf{Lab \& Field Techniques}{: SEM, spectrophotometry, plasmid isolation, column chromatography, cell culture, gene extraction, recombinant protein expression, lithium-ion battery fabrication} \\
     \textbf{Languages}{: English (native), Mandarin Chinese}
    }}
 \end{itemize}


%-------------------------------------------
\end{document}